\chapter{THỰC NGHIỆM VÀ ĐÁNH GIÁ KẾT QUẢ}
\label{chap:chapter4}

\section{Cấu hình thực nghiệm}
\label{sec:cauHinh}

\subsection{Môi trường phần cứng và phần mềm}
\label{subsec:phanCung}

Môi trường thực nghiệm được triển khai trên nền tảng Python với các thư viện
cốt lõi: PyTorch (cài đặt mạng Actor--Critic và thuật toán PPO), Gymnasium
(giao diện môi trường mô phỏng), NumPy và Pandas (xử lý dữ liệu M5 và tính
toán thống kê).

% >>> PLACEHOLDER: Điền cấu hình GPU/CPU cụ thể sau khi hoàn tất thực nghiệm.

\subsection{Nguồn dữ liệu và tiền xử lý}
\label{subsec:duLieu}

Dữ liệu M5 Walmart. Khóa luận sử dụng hai tệp chính từ bộ dữ liệu
M5 \parencite{makridakis2022background, m5dataset}:
\begin{itemize}
    \item \texttt{sales\_train\_evaluation.csv}: chứa $30.490$ chuỗi doanh số
    theo ngày ($3.049$ mặt hàng $\times$ $10$ cửa hàng) trên $1.941$ ngày.
    \item \texttt{calendar.csv}: chứa thông tin lịch, ngày lễ và các sự kiện
    SNAP.
\end{itemize}

Từ tổng số $3.049$ mặt hàng, khóa luận chọn $M = 50$ mặt hàng đại diện cho
$10$ cửa hàng (tương ứng $10$ kho), đảm bảo đa dạng về quy mô nhu cầu: từ các
SKU có nhu cầu thấp ($\lambda < 0{,}5$ đơn vị/ngày) đến các SKU có nhu cầu cao
($\lambda > 5$ đơn vị/ngày). Tham số nhu cầu ($\lambda_i$ trung bình hàng ngày)
cho mỗi cặp kho--mặt hàng được tính từ dữ liệu gốc.

Dữ liệu giả lập (Poisson synthetic). Song song, khóa luận sinh dữ
liệu nhu cầu theo phân phối Poisson $D_t \sim \text{Poisson}(\lambda_i)$ để kiểm
chứng kết quả trong điều kiện kiểm soát.

\subsection{Thiết lập kịch bản mô phỏng}
\label{subsec:kichBan}

Mỗi episode huấn luyện/đánh giá kéo dài $365$ ngày. Bảng~\ref{tab:cauHinhEnv}
tóm tắt cấu hình môi trường.

\begin{table}[htbp]
    \centering
    \caption{Cấu hình môi trường mô phỏng}
    \label{tab:cauHinhEnv}
    \begin{tabular}{|l|c|}
        \hline
        \textbf{Tham số} & \textbf{Giá trị} \\
        \hline
        Số kho ($N$) & $10$ \\ \hline
        Số mặt hàng ($M$) & $50$ \\ \hline
        Số tác tử ($N \times M$) & $500$ \\ \hline
        Độ dài episode & $365$ ngày \\ \hline
        Lead time & $\text{Uniform}(1, 3)$ ngày \\ \hline
        Mức đặt hàng ($K$) & $\{0, 10, 20, 30, 40, 50\}$ \\ \hline
        Ngưỡng fill rate ($\tau$) & $85\%$ \\
        \hline
    \end{tabular}
\end{table}

\section{Phương pháp đối chứng (Baselines)}
\label{sec:baselines}

Hiệu năng của IPPO được so sánh với ba chính sách tồn kho cổ điển:

\begin{itemize}
    \item EOQ \cite{harris1913}: lượng đặt hàng cố định
    $Q^* = \sqrt{2D \cdot c_{dh} / c_{lk}}$, với $D$ tính từ nhu cầu trung bình
    của mỗi SKU. Đặt hàng khi tồn kho giảm xuống dưới ngưỡng $r = \mu_L$.

    \item Chính sách $(s, S)$ \cite{scarf1960}: ngưỡng $s$ và mức
    trần $S$ được tối ưu bằng tìm kiếm lưới (grid search) trên chính môi trường
    mô phỏng, bảo đảm tính công bằng khi so sánh.

    \item Newsvendor \cite{arrow1951}: lượng đặt hàng tối ưu đơn
    kỳ $Q^* = F^{-1}(c_{th}/(c_{th} + c_{lk}))$, áp dụng lại mỗi ngày dựa
    trên phân phối nhu cầu ước lượng.
\end{itemize}

Tất cả baselines được chạy trên cùng một môi trường và cùng
các hạt giống ngẫu nhiên với IPPO để bảo đảm tính so sánh.

\section{Giao thức đánh giá}
\label{sec:giaoThuc}

\subsection{Đánh giá đa hạt giống}
\label{subsec:daHatGiong}

Theo khuyến nghị của Henderson và cộng sự \cite{henderson2018} và Agarwal và cộng sự \cite{agarwal2021}, mọi
kết quả được báo cáo trên tối thiểu $5$ hạt giống ngẫu nhiên. Mỗi seed tương
ứng với:
\begin{itemize}
    \item Khởi tạo trọng số mạng khác nhau.
    \item Chuỗi nhu cầu ngẫu nhiên khác nhau (cùng phân phối).
    \item Chuỗi lead time ngẫu nhiên khác nhau.
\end{itemize}

Kết quả báo cáo gồm giá trị trung bình $\pm$ độ lệch chuẩn. Sự khác biệt giữa
các phương pháp được xác nhận bằng kiểm định thống kê phù hợp (t-test hoặc
Wilcoxon signed-rank test).

\subsection{Các chỉ số đánh giá}
\label{subsec:chiSo}

Bốn nhóm chỉ số chính:

\begin{enumerate}
    \item Tổng chi phí vận hành (Total Cost): tổng $C_t$ theo công
    thức (\ref{eq:tongchiphi}) trên toàn bộ $365$ ngày và $500$ cặp kho--mặt
    hàng.

    \item Tỷ lệ đáp ứng nhu cầu (Fill Rate): $\text{FR} =
    \sum \text{fulfilled} / \sum \text{demand}$, tính trung bình trên toàn hệ
    thống. Mục tiêu $\ge 85\%$.

    \item Cơ cấu chi phí thành phần: phân tách tổng chi phí thành
    holding, stockout, ordering và overflow để phân tích chiến lược mà mỗi
    phương pháp áp dụng.

    \item Đường cong học (Learning Curve): phần thưởng tích lũy trung
    bình qua các episode huấn luyện, chứng minh tính hội tụ.
\end{enumerate}

\section{Kết quả thực nghiệm}
\label{sec:ketQua}

% >>> LƯU Ý: Toàn bộ phần kết quả dưới đây là KHUNG (placeholder).
% >>> CHỈ ĐIỀN SỐ LIỆU SAU KHI ĐÃ CHẠY XONG THỰC NGHIỆM ĐA HẠT GIỐNG
% >>> VÀ CÓ KIỂM ĐỊNH THỐNG KÊ. KHÔNG VIẾT TRƯỚC.

\subsection{Quá trình huấn luyện và tính hội tụ}
\label{subsec:hoiTu}

Đường cong phần thưởng tích lũy trung bình qua các episode huấn luyện cho thấy
mô hình IPPO hội tụ sau khoảng \textit{[cần điền]} episode. Phần thưởng ổn định
và không sụp đổ sau hội tụ, chứng minh tính ổn định của thuật toán PPO với cơ
chế clipped objective.

% >>> Chèn hình đường cong học ở đây khi có:
% \begin{figure}[htbp]
%     \centering
%     \includegraphics[width=0.8\textwidth]{media/learning_curve.png}
%     \caption{Đường cong phần thưởng tích lũy qua các episode huấn luyện}
%     \label{fig:learningCurve}
% \end{figure}

\subsection{So sánh tổng chi phí vận hành}
\label{subsec:soSanhChiPhi}

Bảng~\ref{tab:ketQuaChiPhi} trình bày tổng chi phí vận hành sau $365$ ngày,
trung bình trên $5$ hạt giống.

\begin{table}[htbp]
    \centering
    \caption{Tổng chi phí vận hành (trung bình $\pm$ độ lệch chuẩn, $5$ seeds)}
    \label{tab:ketQuaChiPhi}
    \begin{tabular}{|l|c|c|c|}
        \hline
        \textbf{Phương pháp} & \textbf{Total Cost} & \textbf{Fill Rate} &
        \textbf{$p$-value} \\
        \hline
        EOQ           & \textit{[cần điền]} & \textit{[cần điền]} & --- \\ \hline
        $(s, S)$      & \textit{[cần điền]} & \textit{[cần điền]} & --- \\ \hline
        Newsvendor    & \textit{[cần điền]} & \textit{[cần điền]} & --- \\ \hline
        \textbf{IPPO} & \textit{[cần điền]} & \textit{[cần điền]} &
        \textit{[cần điền]} \\
        \hline
    \end{tabular}
\end{table}

\subsection{Phân tích cơ cấu chi phí}
\label{subsec:coCauChiPhi}

% >>> Chèn bảng và/hoặc biểu đồ cột so sánh holding/stockout/ordering/overflow
% >>> giữa 4 phương pháp khi có số liệu.

Phân tích cơ cấu chi phí cho thấy IPPO \textit{[cần phân tích: chiến lược
đặt hàng thiên về tần suất thấp lô lớn hay tần suất cao lô nhỏ so với
baselines?]}.

\subsection{Đánh giá tỷ lệ đáp ứng dịch vụ}
\label{subsec:danhGiaFillRate}

Tỷ lệ đáp ứng trung bình của IPPO đạt \textit{[cần điền]}\%, so với ngưỡng
mục tiêu $85\%$. Phân tích theo nhóm SKU (nhu cầu cao/thấp/gián đoạn) cho thấy
\textit{[cần phân tích: IPPO có duy trì fill rate đồng đều hay có nhóm SKU bị
hy sinh?]}.

\section{Phân tích chuyên sâu}
\label{sec:phanTich}

\subsection{Khả năng tổng quát hóa qua parameter sharing (Câu hỏi 2)}
\label{subsec:tongQuatHoa}

Để trả lời Câu hỏi nghiên cứu 2, khóa luận phân tích hiệu năng của chính sách
chia sẻ tham số theo nhóm SKU có quy mô nhu cầu khác nhau. Các SKU được phân
thành $3$ nhóm dựa trên $\lambda_i$ trung bình: nhu cầu thấp ($\lambda < 1$),
trung bình ($1 \le \lambda < 5$), và cao ($\lambda \ge 5$).

% >>> Điền bảng/biểu đồ fill rate và cost theo nhóm SKU khi có số liệu.

\subsection{Ảnh hưởng của hàm phần thưởng phân rã cục bộ (Câu hỏi 3)}
\label{subsec:ablationReward}

Để trả lời Câu hỏi nghiên cứu 3, khóa luận so sánh hai cấu hình:
\begin{itemize}
    \item Local reward (công thức \ref{eq:localReward}): mỗi tác tử
    nhận phần thưởng riêng.
    \item Global reward: tất cả tác tử nhận cùng phần thưởng
    $r_{\text{global}} = \sum_i r_i / (N \times M)$.
\end{itemize}

% >>> Điền kết quả ablation study khi có số liệu.

\subsection{Các kịch bản IPPO chưa vượt trội}
\label{subsec:hanChe}

Khóa luận trình bày trung thực những kịch bản mà IPPO không đạt hiệu năng tốt
hơn chính sách cổ điển, kèm phân tích nguyên nhân:

% >>> Điền phân tích khi có kết quả thực nghiệm cụ thể.
% >>> Ví dụ: nhu cầu hoàn toàn ổn định (Poisson thấp) thì (s,S) đã gần tối ưu;
% >>> hoặc SKU rất hiếm (lumpy demand) khiến tín hiệu học quá thưa.

\section{Thảo luận}
\label{sec:thaoLuan}

% >>> Tổng hợp và thảo luận kết quả, trả lời 3 câu hỏi nghiên cứu.
% >>> So sánh với kết quả của Oroojlooyjadid (2022) và Gijsbrechts (2022).
% >>> Nhận xét về tính khả thi áp dụng trong thực tế.

Kết quả thực nghiệm cho thấy \textit{[cần tổng hợp sau khi có đầy đủ số liệu]}.